\chapter{wuffle}\label{ch:wuffle}
\begin{goals}
\textbf{In this chapter} you learn what wuffle is: one file with two homes, a project that keeps each
dataset once and shows it through layers, and a posture that keeps your data on your machine.
\end{goals}

wuffle is a stereonet program in a single HTML file. It has no installer and no server, and it needs no
account. You open the file in a browser and it runs there, offline.

The name puns on the Wulff net. Orientation data belongs on an \term[equal-area net]{equal-area}
(Schmidt) net, though, so that is what wuffle draws unless you ask for another projection.

\section{Two homes, one file}\label{sec:two-homes}

The same file is two programs, chosen by where it opens:

\begin{description}
\item[The field.] On a phone, wuffle is a \term[compass]{compass--clinometer}. Lay the phone on a
bedding plane and it reads dip direction and dip; point it along a lineation and it reads trend and
plunge. Each measurement lands on a live net and goes into a log you can publish
(\cref{ch:field}).
\item[The desk.] On a computer, wuffle is a \term{workbench}: nets, rose diagrams and fabric plots,
selections drawn on the sphere, statistics, a shell, and figures for a report.
\end{description}

The two meet in the project. A log published from the phone opens on the desk with its signature
checked, and its original bytes are kept untouched beside the data read from them
(\cref{sec:field-logs}).

\section{Data once, layers per diagram}\label{sec:data-once}

A wuffle project holds three kinds of thing:

\begin{description}
\item[Datasets.] A \term{dataset} is a table of measurements: rows, columns, and which columns hold the
attitude. It lives in the project once, whatever draws it.
\item[Diagrams.] A \term{diagram} is a net, a rose diagram or a fabric plot. Each has its own settings
(projection, hemisphere, rotation) and its own tab.
\item[Layers.] A \term{layer} is one dataset drawn on one diagram, with its own style: poles or great
circles, contours, a mean, eigenvectors.
\end{description}

Every layer of a dataset is a view of the same rows. Edit a row and every view changes. Select poles
on one net and the same rows are selected on every other net and in the table, because the selection
belongs to the data, not to the drawing.

\begin{why}
A dataset copied into each diagram goes out of step the first time one copy is edited, and then you
can no longer tell which net shows the data you mean. One dataset with many views can't drift: there is
only one set of rows to be right or wrong.
\end{why}

Results are datasets too. A mean, the eigenvectors, a $\beta$ intersection, or a selection made into a
set becomes a new dataset that records where it came from, its \term{lineage}, so it can be drawn,
styled, saved and questioned like the measurements it was computed from.

\section{Local, and honest about it}\label{sec:local}

wuffle fetches nothing when it opens. Your data stays in the browser tab, and between saves wuffle keeps
a copy of the project in the browser's own storage, on your machine (\cref{sec:autosave}). The only
things that can reach the network are things you type in the console, such as a package install.

The badge in the corner says \ui{● local}. It does not say \emph{sealed}: a sealed page is one whose
browser policy forbids the network outright, and wuffle doesn't carry that policy, so it makes the
smaller claim.

\section{What it is built on}

wuffle's maths is \code{@gcu/bearing}, a library of orientation geometry: projections, the Fisher and
eigen statistics, contouring, rotations. The nets, selections and figure frames come from
\code{@gcu/stereonet}, which the GCU notebook's stereonet cell also uses. The console is
\code{geas}, a shell written for the browser. All of it ships inside the one file.
